\documentclass[11pt]{article}

\usepackage[margin=1in]{geometry}
\usepackage{amsmath,amssymb}
\usepackage{booktabs}
\usepackage{tabularx}
\usepackage{array}
\usepackage{enumitem}
\usepackage[T1]{fontenc}
\usepackage{hyperref}

\newcommand{\code}[1]{\texttt{#1}}
\newcolumntype{L}{>{\raggedright\arraybackslash}X}
\emergencystretch=3em

\title{Orientational-order diagnostics in \code{cluster\_diagnostics.py}\\
and \code{cluster\_diagnostics\_3d.py}: definitions and interpretation}
\author{FHDeX}
\date{}

\begin{document}
\maketitle

\noindent This note explains the orientational-order diagnostics written by
\code{cluster\_diagnostics.py} (2D) and \code{cluster\_diagnostics\_3d.py} (3D):
what each quantity is and how to read it. The Markdown version is
\code{README\_order\_diagnostics.md}.

\section{\texorpdfstring{Why $g(r)$ is not enough}{Why g(r) is not enough}}

The radial distribution function $g(r)$ counts pairs of clusters at distance
$r$, averaged over every direction. It records how far apart clusters are, but
not in which directions their neighbours lie. A crystal and a dense amorphous
packing with the same spacing therefore have $g(r)$ peaks in the same places;
the crystal's peaks are only somewhat sharper. That is the case in
\code{exec/dean\_kow/interaction}:
\begin{itemize}[nosep]
  \item the 2D run \code{\_new\_repul\_2} looks like a hexagonal lattice;
  \item the 3D run \code{\_3d\_noise} looks amorphous;
  \item the cluster $g(r)$ of both has the same nearest-neighbour peak.
\end{itemize}
Telling these states apart needs angular information. The tools measure it in
two independent ways:
\begin{enumerate}[nosep]
  \item \textbf{Bond-orientational order in real space:} the directions from
        each cluster centre to its neighbours, compared within each
        neighbourhood (local order) and across the box (global order).
  \item \textbf{Bragg spots in $k$ space:} whether the main peak of the
        structure factor $S(\mathbf k)$ sits on a few wave vectors (spots) or
        is spread around the ring.
\end{enumerate}
Both are written for every frame to \code{order.txt}. The per-cluster bond
order is added as extra columns of \code{clusters.txt}.

\section{Neighbours and bonds}

Order is measured on the cluster centres $\mathbf x_i$ (the circular-mean
centres of mass in \code{clusters.txt}). The bond vectors are
$\mathbf r_{ij} = \mathbf x_j - \mathbf x_i$, with the minimum image in each
periodic direction. The neighbours of $i$ are
\begin{itemize}[nosep]
  \item by default, its \code{-{}-nn} nearest centres (6 in 2D, 12 in 3D);
  \item with \code{-{}-nn-cut R}, every centre with $|\mathbf r_{ij}|<R$. A
        good $R$ is the first minimum of the cluster $g(r)$ after its first
        peak. The number of neighbours $n_i$ then varies; a cluster with no
        neighbours gets NaN and is left out of the averages.
\end{itemize}
The neighbour count should match the coordination of the packing being tested
for: 6 for 2D hexagonal, 12 for fcc or hcp, 8 (or 14 with the second shell) for
bcc, and 6 for simple cubic. Bonds are not symmetrized: $j$ can be among $i$'s
nearest neighbours without $i$ being among $j$'s. With a fixed \code{-{}-nn}
this is the usual convention and has little effect.

\section{\texorpdfstring{Bond order in 2D: $\psi_6$}{Bond order in 2D: psi6}}

With $\theta_{ij}$ the angle of $\mathbf r_{ij}$ to the $x$ axis, each cluster
gets the complex number
\begin{equation}
  \psi_{6,i} = \frac{1}{n_i}\sum_{j} e^{6i\theta_{ij}} .
\end{equation}
The factor 6 makes the six bonds of a perfect hexagon contribute the same
phase, since rotating a bond by $60^\circ$ changes $6\theta$ by $360^\circ$.
The modulus $|\psi_{6,i}|$ says how hexagonal the neighbourhood is (1 for a
perfect hexagon); the phase $\arg\psi_{6,i}/6$ is the hexagon's orientation.
Over the $N$ clusters of a frame,
\begin{equation}
  \psi_6^{\mathrm{loc}} = \frac1N\sum_i|\psi_{6,i}|
  \quad(\code{psi6\_local}),\qquad
  \psi_6^{\mathrm{glob}} = \Bigl|\frac1N\sum_i\psi_{6,i}\Bigr|
  \quad(\code{psi6\_global}).
\end{equation}
Taking the modulus first ignores orientation; averaging the complex numbers
first lets differently oriented hexagons cancel. Always
$\psi_6^{\mathrm{glob}}\le\psi_6^{\mathrm{loc}}$. \code{clusters.txt} gets a
\code{psi6} column with $|\psi_{6,i}|$.

\begin{center}
\begin{tabularx}{\textwidth}{@{}L l L@{}}
\toprule
Configuration (6 nearest neighbours) & $\psi_6^{\mathrm{loc}}$ & $\psi_6^{\mathrm{glob}}$ \\
\midrule
perfect hexagonal lattice & 1 & 1 \\
hexagonal lattice, random displacements of 5\% of the spacing & 0.92 & 0.92 \\
random points (Poisson) & 0.36 & $\approx0.4/\sqrt N$ (0.054 for $N=80$, 0.018 for $N=500$) \\
\bottomrule
\end{tabularx}
\end{center}
The random local value is the mean modulus of an average of six random
phases, $\approx\sqrt{\pi/24}\approx0.36$. The global value of random
orientations is not zero in a finite system, so it must be compared with
$0.4/\sqrt N$, not with 0.

\section{\texorpdfstring{Bond order in 3D: Steinhardt $q_4$, $q_6$}{Bond order in 3D: Steinhardt q4, q6}}

In 3D a bond has a direction $\hat{\mathbf r}_{ij}$ rather than one angle. Each
neighbourhood is expanded in spherical harmonics:
\begin{equation}
  q_{lm,i} = \frac1{n_i}\sum_j Y_{lm}(\hat{\mathbf r}_{ij}),\quad m=-l,\dots,l,
  \qquad
  q_{l,i} = \Bigl(\frac{4\pi}{2l+1}\sum_{m=-l}^{l}|q_{lm,i}|^2\Bigr)^{1/2}.
\end{equation}
These are the Steinhardt--Nelson--Ronchetti bond-order parameters.
\begin{itemize}[nosep]
  \item The $2l+1$ complex numbers $q_{lm,i}$ are the 3D analogue of
        $\psi_{6,i}$: they encode both the shape and the orientation of the
        neighbourhood.
  \item The sum over $m$ makes $q_{l,i}$ rotationally invariant, like
        $|\psi_6|$ in 2D. The normalization gives $q_l=1$ when all bonds
        point in one direction.
  \item $l=4$ and $l=6$ are the standard pair: cubic and icosahedral
        symmetries have strong $l=4$ and $l=6$ components.
  \item $Y_{lm}$ is computed from Legendre polynomials in
        \code{numpy.polynomial}, without scipy.
\end{itemize}
The averages follow the 2D logic:
\begin{equation}
  q_l^{\mathrm{loc}} = \frac1N\sum_i q_{l,i}
  \quad(\code{q4},\ \code{q6}),\qquad
  Q_l^{\mathrm{glob}} = \Bigl(\frac{4\pi}{2l+1}\sum_m
     \Bigl|\frac1N\sum_i q_{lm,i}\Bigr|^2\Bigr)^{1/2}
  \quad(\code{Q4\_global},\ \code{Q6\_global}).
\end{equation}
The global value equals the local value for a single perfect crystal of any
orientation and is near zero when orientations are random.
\code{clusters.txt} gets \code{q4} and \code{q6} columns.

\subsection*{Crystal packings}
The reference lattices are the standard ways of packing equal spheres (here,
cluster centres) in 3D:
\begin{center}
\begin{tabularx}{\textwidth}{@{}l L l l@{}}
\toprule
Name & Structure & Neighbours & Packing fraction \\
\midrule
fcc, face-centred cubic & a cube with a point at each corner and at the centre of each face (4 points per cell) & 12 & 0.74 (densest) \\
hcp, hexagonal close-packed & hexagonal layers stacked ABAB\dots & 12 & 0.74 \\
bcc, body-centred cubic & a cube with a point at each corner and one at its centre (2 per cell) & 8 (+6 at 1.15$\times$) & 0.68 \\
sc, simple cubic & points only at the cube corners (1 per cell) & 6 & 0.52 \\
\bottomrule
\end{tabularx}
\end{center}
\begin{itemize}
  \item \textbf{fcc and hcp} are stacks of close-packed layers, each identical
        to the 2D hexagonal lattice. Each layer sits in the hollows of the one
        below, which allows three positions A, B, C: ABAB\dots\ stacking gives
        hcp and ABCABC\dots\ gives fcc. Random stackings are common in
        colloidal crystals. Both have 12 equidistant neighbours, so $g(r)$
        barely separates them; their angular arrangements, and so $q_4$ and
        $q_6$, differ.
  \item \textbf{bcc} is less dense and more open. Each point is at the centre
        of a cube of 8 neighbours, with 6 more at $2/\sqrt3\approx1.15$ times
        that distance. Its reference values use 8 neighbours. With 12, the
        first and second shells are mixed, so use \code{-{}-nn 8} (or 14, or
        \code{-{}-nn-cut} between the shells) when testing for bcc.
  \item \textbf{sc} is rarely stable for equal spheres. It is listed because
        its $q_4$ and $q_6$ are distinctive.
  \item \textbf{Which forms:}
        \begin{itemize}[nosep]
          \item hard or short-range repulsion typically gives fcc or hcp;
          \item soft, long-range repulsion often gives bcc, especially near
                melting;
          \item soft, penetrable potentials, such as the gema and HK
                potentials of the interaction code, can give fcc or bcc
                cluster crystals depending on density and noise.
        \end{itemize}
  \item A \textbf{polycrystal} is many crystalline grains with different
        orientations. An \textbf{amorphous} (glassy or liquid-like) packing
        has no lattice, only short-range order.
\end{itemize}

\subsection*{Reference values}
\begin{center}
\begin{tabular}{@{}l l l@{}}
\toprule
Packing (neighbours) & $q_4$ & $q_6$ \\
\midrule
fcc (12) & 0.191 & 0.575 \\
hcp (12) & 0.097 & 0.485 \\
bcc (8) & 0.509 & 0.629 \\
simple cubic (6) & 0.764 & 0.354 \\
icosahedral cluster (12), for comparison & 0 & 0.663 \\
random points (12) & 0.28 & 0.28; global $Q_6\approx0.4/\sqrt N$ (0.017 for $N=500$) \\
\bottomrule
\end{tabular}
\end{center}
The lattice values were checked with the code: they do not depend on the
lattice's orientation, and global $=$ local. Thermal displacements lower all
crystal values and broaden their distributions. A dense liquid or glass
typically has $q_6\approx0.35$--$0.45$: above random, because neighbours keep a
minimum distance, but below the crystals.

\section{\texorpdfstring{Bragg spots or a ring: $S(\mathbf k)$ on the peak shell}{Bragg spots or a ring}}

For each frame the tool computes the full structure factor
$S(\mathbf k)=|\hat n(\mathbf k)|^2/N_{\mathrm{tot}}$ on the FFT grid, before
any shell average, and takes the shell (a ring in 2D, a spherical shell in 3D)
with the largest mean $S$. That is normally the main peak at
$2\pi/(\text{cluster spacing})$. On that shell, with $n$ wave vectors and values
$S_1,\dots,S_n$,
\begin{equation}
  \mathrm{PR} = \frac{\bigl(\sum_a S_a\bigr)^2}{n\sum_a S_a^2},\qquad
  \code{top\_share} = \frac{\text{sum of the \code{-{}-top-modes} largest } S_a}{\sum_a S_a}.
\end{equation}
\code{order.txt} also gives \code{kpeak}, the mean $|\mathbf k|$ of the shell,
and \code{nmodes} $=n$. PR is a participation ratio: if $S$ were equal on $k$
of the $n$ modes and zero elsewhere, $\mathrm{PR}=k/n$ exactly. It assumes no
particular lattice.
\begin{itemize}[nosep]
  \item \textbf{Crystal:} a few Bragg spots carry the shell, so
        $\mathrm{PR}\approx(\text{number of spots})/n$ and
        $\code{top\_share}\approx1$ when \code{-{}-top-modes} is at least the
        number of spots. Spots come in $\pm\mathbf k$ pairs: a 2D hexagonal
        lattice has 6 on its first ring, fcc 8 on its first shell, bcc 12.
  \item \textbf{Isotropic ring} (amorphous, liquid): with random phases each
        $S_a$ is roughly exponentially distributed, so
        $\langle S^2\rangle=2\langle S\rangle^2$ and $\mathrm{PR}\approx0.5$;
        \code{top\_share} is about $\code{top\_modes}/n$, somewhat more due to
        fluctuations.
  \item \textbf{Uniform field} (e.g.\ $t=0$): PR and \code{top\_share} are NaN.
\end{itemize}

\begin{center}
\begin{tabular}{@{}l r l l@{}}
\toprule
Field & $n$ & PR & \code{top\_share} \\
\midrule
2D hexagonal lattice of Gaussian blobs & 60 & $0.100=6/60$ & 1.00 \\
2D random blobs & 12 & 0.61 & --- \\
3D fcc lattice of blobs, $48^3$ & 602 & $0.0133=8/602$ & 1.00 \\
3D random blobs & 98 & 0.55 & 0.36 \\
\bottomrule
\end{tabular}
\end{center}
This measure is independent of the cluster detection (\code{-{}-threshold},
\code{-{}-mmin}) and of the neighbour rule, so it cross-checks the bond order.

\paragraph{Seeing the spots.} \code{-{}-sf-full} writes the full
$S(\mathbf k)$, averaged over the frames in \code{-{}-t-range}, as the plotfile
\code{sf\_full\_avg} ($k=0$ at the centre, components \code{S} and
\code{log10S}); \code{-{}-sf-full-per-frame} writes one per frame. In VisIt,
ParaView or amrvis, a 2D crystal shows distinct spots (six on the first ring
for a hexagonal lattice, more at $\sqrt3$ and 2 times that radius), a
polycrystal several sets of spots on one ring, and an amorphous state a
uniform speckled ring. In 3D use slices through $k=0$ or an isosurface of
\code{log10S}: crystal spots are about $n/(\text{spots})\gtrsim100$ times the
shell mean, whereas random speckle on $n$ modes reaches only about $\ln n$
times the mean ($\approx7$ for 1000 modes).


\section{\texorpdfstring{The columns of \code{order.txt}}{The columns of order.txt}}
One row per plotfile, after a header line giving the neighbour rule and
\code{-{}-top-modes}.

\subsection*{2D (\code{cluster\_diagnostics.py})}
\begin{center}
\begin{tabularx}{\textwidth}{@{}r l L@{}}
\toprule
\# & Column & Meaning \\
\midrule
1 & \code{t} & plotfile time \\
2 & \code{step} & plotfile step \\
3 & \code{nclusters} & number of clusters in the frame (as in \code{summary.txt}) \\
4 & \code{psi6\_local} & $\psi_6^{\mathrm{loc}}=\overline{|\psi_{6,i}|}$: how hexagonal each neighbourhood is, whatever its orientation; 1 for a perfect hexagon, $\approx0.36$ for random points \\
5 & \code{psi6\_global} & $\psi_6^{\mathrm{glob}}=|\overline{\psi_{6,i}}|$: one orientation across the box; $\approx\psi_6^{\mathrm{loc}}$ for a single crystal, $\approx0.4/\sqrt N$ with no common orientation \\
6 & \code{kpeak} & mean $|\mathbf k|$ of the $S(k)$ ring with the largest mean $S$ in this frame, normally the main peak at $\approx2\pi/\text{spacing}$ \\
7 & \code{nmodes} & number of FFT wave vectors $n$ on that ring \\
8 & \code{PR} & $(\sum S)^2/(n\sum S^2)$, roughly the fraction of modes carrying the signal: $\approx$ spots$/n$ for Bragg spots, $\approx0.5$ for a uniform ring \\
9 & \code{top\_share} & fraction of the ring's $S$ in its \code{-{}-top-modes} largest modes (default 6): $\approx1$ for a crystal, small for a ring \\
\bottomrule
\end{tabularx}
\end{center}
Example, the last frame of \code{\_new\_repul\_2}:
\begin{verbatim}
0.08773803711 400000 53 0.987506 0.986861 50.306732 48 0.125774 0.996647
\end{verbatim}
At $t=0.0877$ (step 400000) there are 53 clusters. Local and global $\psi_6$
are both 0.99: one hexagonal crystal with a single orientation. The $S(k)$
peak ring at $|\mathbf k|\approx50.3$ has 48 modes, and $\mathrm{PR}=0.126\approx6/48$,
with the 6 largest modes holding 99.7\% of it: six Bragg spots. The $t=0$ row,
\code{0 0 0 nan nan 7.58 8 nan nan}, is a uniform field: the order values are
NaN, and \code{kpeak} and \code{nmodes} only describe the first ring.

\subsection*{3D (\code{cluster\_diagnostics\_3d.py})}
\begin{center}
\begin{tabularx}{\textwidth}{@{}r l L@{}}
\toprule
\# & Column & Meaning \\
\midrule
1--3 & \code{t}, \code{step}, \code{nclusters} & as in 2D \\
4 & \code{q4} & $\overline{q_{4,i}}$, local (fcc 0.19, bcc 0.51 with 8 nn, random 0.28) \\
5 & \code{q6} & $\overline{q_{6,i}}$, local (fcc 0.575, hcp 0.485, random 0.28) \\
6 & \code{Q4\_global} & $Q_4$ from $q_{4m}$ averaged over all clusters \\
7 & \code{Q6\_global} & $Q_6$ from $q_{6m}$ averaged over all clusters: $\approx q_6$ for a single crystal, $\approx0.4/\sqrt N$ with no common orientation \\
8--11 & \code{kpeak}, \code{nmodes}, \code{PR}, \code{top\_share} & as in 2D; \code{-{}-top-modes} defaults to 12 \\
\bottomrule
\end{tabularx}
\end{center}
Example, the last frame of \code{\_3d\_noise}:
\begin{verbatim}
0.03797721181 60000 476 0.126799 0.424355 0.016313 0.073719
    56.909256 1142 0.318627 0.096069
\end{verbatim}
(one line in the file). Local $q_6=0.42$, $q_4=0.13$ lie between random and the
crystals and match no crystal; global $Q_6=0.074$ is six times below local, so
there is no common orientation; on the 1142-mode ring at
$|\mathbf k|\approx56.9$, $\mathrm{PR}=0.32$ and the 12 largest modes hold only
9.6\%: a ring, not spots. This is an amorphous packing.

\section{How to read them together}

\begin{center}
\begin{tabularx}{\textwidth}{@{}L L L L@{}}
\toprule
Local order & Global order & $S(k)$ peak shell & Interpretation \\
\midrule
near the crystal value & $\approx$ local & PR $\approx$ spots$/n$, \code{top\_share} $\approx1$ & \textbf{single crystal} spanning the box \\
near the crystal value & well below local, well above $0.4/\sqrt N$ & intermediate PR, several groups of spots & \textbf{polycrystal} of a few grains; global $\approx$ local$/\sqrt{\text{grains}}$ \\
near the crystal value & $\approx0.4/\sqrt N$ & ring, PR $\approx0.3$--$0.5$ & \textbf{fine polycrystal}, hard to tell from amorphous; use per-cluster histograms \\
between random and crystal & $\approx0.4/\sqrt N$ & ring & \textbf{amorphous / dense liquid} \\
$\approx$ random & $\approx0.4/\sqrt N$ & ring or no clear peak & \textbf{disordered}, gas-like \\
\bottomrule
\end{tabularx}
\end{center}
In 2D, dense liquids and hexatic phases can also have fairly high local
$\psi_6$, so local $\psi_6$ alone does not prove a crystal: the global value
and the spots decide.

\paragraph{Time series.} Plot the \code{order.txt} columns against $t$.
\begin{itemize}[nosep]
  \item \emph{Crystallization:} local order rises first, then global order,
        while PR drops towards spots$/n$ and \code{top\_share} rises to 1.
  \item \emph{Coarsening of a polycrystal:} local order flat, global order
        rising slowly as grains merge.
  \item \emph{Amorphous or glassy:} everything roughly constant. A slow drift
        may mean very slow ordering; run longer before concluding.
\end{itemize}

\paragraph{Per-cluster distributions.} The \code{clusters.txt} columns
(\code{psi6}, or \code{q4 q6}) give the distribution over clusters.
\begin{itemize}[nosep]
  \item One narrow peak at a crystal value: crystal.
  \item One broad peak below the crystal values: amorphous.
  \item Two peaks: crystalline grains with disordered regions; the weight of
        the upper peak estimates the crystalline fraction.
  \item 3D, $q_4$ against $q_6$: points near $(0.19, 0.575)$ for fcc,
        $(0.10, 0.485)$ for hcp, $(0.51, 0.63)$ for bcc (8 neighbours).
\end{itemize}

\section{\texorpdfstring{Worked example: \code{\_new\_repul\_2} (2D) and \code{\_3d\_noise} (3D)}{Worked example}}

\begin{center}
\begin{tabularx}{\textwidth}{@{}l L L@{}}
\toprule
 & 2D \code{\_new\_repul\_2} & 3D \code{\_3d\_noise} \\
\midrule
frames & $t=0.022$ to $0.088$ & $t=0.032$ to $0.038$ \\
clusters $N$ & 41--53 & 465--480 \\
local order & $\psi_6$: $0.63\to0.89\to0.91\to0.988$ & $q_4=0.13$, $q_6=0.41$--$0.42$ \\
global order & $\psi_6$: $0.55\to0.86\to0.89\to0.987$ & $Q_4=0.012$--$0.016$, $Q_6=0.058$--$0.074$ \\
noise floor $0.4/\sqrt N$ & 0.055 & 0.018 \\
$S(k)$ peak shell & 48 modes; PR $0.23\to0.126$; \code{top\_share} (6) $0.66\to0.997$ & 1142 modes; PR 0.32; \code{top\_share} (12) 0.08--0.10 \\
\bottomrule
\end{tabularx}
\end{center}

\paragraph{2D: a single hexagonal crystal, formed during the run.} Local and
global $\psi_6$ agree and approach 1. Six Bragg spots carry essentially the
whole ring: PR $=0.126$ against $6/48=0.125$. Early on the ordering is
incomplete ($\psi_6=0.63$, PR $=0.23$); the lattice anneals by $t\approx0.09$.

\paragraph{3D: amorphous.} Local $q_6=0.42$ is above random (0.28) but below
hcp (0.485) and fcc (0.575), and $q_4=0.13$ matches no crystal: a dense,
liquid-like packing. Global $Q_6=0.07$ is about 6 times smaller than the local
value, so neighbourhoods share no orientation; it is only 4 times the
random-orientation floor, so at most there are very many tiny grains. The
$S(k)$ peak is a ring: the 12 largest of 1142 modes hold under 10\% of it, and
PR $=0.32$ is far above the $\approx0.01$ Bragg spots would give. In
\code{sf\_full\_avg} over $t=0.037$--$0.038$ the brightest modes on the peak
shell are about 14 times the shell mean: above pure speckle ($\approx7$) but
far below Bragg spots ($\gtrsim100$), so the ring is only weakly lumpy. The values
drift only slowly over the window, so slow ordering at later times cannot be
excluded.

\section{Caveats}
\begin{itemize}
  \item \textbf{Few clusters:} the global floor is $\approx0.4/\sqrt N$; with
        10--20 clusters, global values of 0.1--0.15 can arise by chance.
  \item \textbf{Neighbour rule:} use a \code{-{}-nn} matching the expected
        coordination (\code{-{}-nn 8} for bcc), or \code{-{}-nn-cut} at the
        first minimum of the cluster $g(r)$ for loose or polydisperse
        packings. Different rules shift the reference values; compare
        against values computed with the same rule.
  \item \textbf{Cluster detection:} stray small clusters inside a lattice, or a
        merger in progress, lower the local order. Check \code{-{}-mmin} and
        \code{-{}-threshold} if it is unexpectedly low. The $S(k)$ measure
        does not depend on this.
  \item \textbf{Periodic box:} a crystal fits the box without strain only for
        certain sizes and orientations. One that does not fit has defects or a
        twist, which lowers global order and smears the spots. The clean spots
        of \code{\_new\_repul\_2} show that its lattice fits the box.
  \item \textbf{Peak shell:} the shell of largest mean $S$ is used. If a low-$k$
        shell is larger, for example during coarsening, \code{kpeak} jumps and
        PR describes that shell instead; check \code{kpeak} against
        $2\pi/\text{spacing}$.
  \item \textbf{Not implemented:} classifying individual clusters as
        crystalline with the ten Wolde criterion (normalized
        $\mathbf q_{6,i}\cdot\mathbf q_{6,j}>0.5$ for most neighbours), which
        would give crystalline fractions and grain sizes directly.
\end{itemize}

\section*{References}
\begin{itemize}[nosep]
  \item P.~J. Steinhardt, D.~R. Nelson and M. Ronchetti, Phys.\ Rev.\ B
        \textbf{28}, 784 (1983): bond-orientational order, $q_4$ and $q_6$.
  \item D.~R. Nelson and B.~I. Halperin, Phys.\ Rev.\ B \textbf{19}, 2457
        (1979): $\psi_6$ and 2D melting.
  \item P.~R. ten Wolde, M.~J. Ruiz-Montero and D. Frenkel, J.\ Chem.\ Phys.\
        \textbf{104}, 9932 (1996): per-particle crystallinity from $q_6$
        correlations.
  \item W. Lechner and C. Dellago, J.\ Chem.\ Phys.\ \textbf{129}, 114707
        (2008): averaged $q_4$, $q_6$ for better separation of crystal types.
\end{itemize}

\end{document}
